% ============================================================================
% LEHRERHINWEISE V2: Mi instituto - Objetos de clase & hay
% Klasse 7 | Spanisch | Version 2
% ============================================================================

\documentclass[11pt,a4paper]{article}

\usepackage[utf8]{inputenc}
\usepackage[T1]{fontenc}
\usepackage[ngerman]{babel}
\usepackage{geometry}
\usepackage{fancyhdr}
\usepackage{xcolor}
\usepackage{enumitem}
\usepackage{tabularx}
\usepackage{array}
\usepackage{tcolorbox}
\usepackage{booktabs}

% ============================================================================
% LAYOUT
% ============================================================================
\geometry{left=2cm, right=2cm, top=2cm, bottom=2cm}
\definecolor{fach}{HTML}{c41e3a}
\definecolor{grau}{HTML}{666666}
\definecolor{hellgrau}{HTML}{f5f5f5}
\definecolor{basis}{HTML}{2a7a4b}
\definecolor{standard}{HTML}{2c5aa0}
\definecolor{erweiterung}{HTML}{7b1fa2}

\pagestyle{fancy}
\fancyhf{}
\fancyhead[L]{\small\textcolor{grau}{Spanisch Klasse 7 -- Unidad 2}}
\fancyhead[R]{\small\textcolor{grau}{LH-01-V2 (Lehrerhinweise)}}
\fancyfoot[C]{\small\thepage}
\renewcommand{\headrulewidth}{0.4pt}

% ============================================================================
% DOKUMENT
% ============================================================================
\begin{document}

% ----------------------------------------------------------------------------
% KOPFBEREICH
% ----------------------------------------------------------------------------
\begin{center}
{\Large\textcolor{fach}{\textbf{Lehrerhinweise}}}\\[0.3cm]
{\large Objetos de clase -- hay}\\[0.2cm]
{\normalsize Doppelstunde (90 Min.) | Qué pasa 1, S. 36--38}
\end{center}

\vspace{0.5cm}
\hrule
\vspace{0.5cm}

% ----------------------------------------------------------------------------
% KURZÜBERSICHT
% ----------------------------------------------------------------------------
\section*{Kurzübersicht}

\begin{tabular}{ll}
\textbf{Thema:} & Klassenzimmer-Vokabeln + hay (es gibt) \\
\textbf{Lehrwerk:} & Qué pasa 1, Unidad 2, S. 36--38 \\
\textbf{Materialien:} & AB-01-V2, LP-01-V2 (optional), Bildkarten \\
\textbf{Fokus:} & 8--12 Vokabeln, Struktur hay + un/una/Zahl \\
\textbf{NICHT:} & estar, Lektionstext, -er/-ir Verben \\
\end{tabular}

\vspace{0.5cm}

% ----------------------------------------------------------------------------
% STUNDENVERLAUF KOMPAKT
% ----------------------------------------------------------------------------
\section*{Stundenverlauf (Kompaktplan)}

\subsection*{1. Stunde (45 Min.)}

\begin{tabular}{|p{1.2cm}|p{2.5cm}|p{9cm}|}
\hline
\textbf{Zeit} & \textbf{Phase} & \textbf{Inhalt} \\
\hline
0--3 & Einstieg & Begrüßung, auf Gegenstände zeigen: \textit{¿Qué es esto?} \\
\hline
3--12 & Input & 8 Vokabeln einführen (chorisch + einzeln nachsprechen) \\
\hline
12--18 & Übung I & Zeigen + Benennen: \textit{Es una mesa.} \\
\hline
18--25 & Grammatik & hay induktiv: Sätze an Tafel, Regel ableiten \\
\hline
25--35 & Übung II & Kettenübung: \textit{En mi clase hay...} \\
\hline
35--42 & Sicherung & Hefteintrag (Vokabeln + hay-Regel) \\
\hline
42--45 & Ausblick & Ankündigung: \textit{¿Qué hay en tu mochila?} \\
\hline
\end{tabular}

\subsection*{2. Stunde (45 Min.)}

\begin{tabular}{|p{1.2cm}|p{2.5cm}|p{9cm}|}
\hline
\textbf{Zeit} & \textbf{Phase} & \textbf{Inhalt} \\
\hline
0--5 & Warm-up & Bildkarten-Wettbewerb (schnell benennen) \\
\hline
5--8 & Demo & LK zeigt eigene Tasche: \textit{En mi mochila hay...} \\
\hline
8--18 & PA & Dialog: \textit{¿Qué hay en tu mochila?} \\
\hline
18--22 & Präsentation & Einige Paare stellen vor \\
\hline
22--25 & Erweiterung & \textit{¿Cuántos/as ... hay?} einführen \\
\hline
25--37 & EA & AB-01-V2 bearbeiten (Aufg. 1--5) \\
\hline
37--42 & Sicherung & Lösungen besprechen \\
\hline
42--45 & Abschluss & HA: Vokabeln, optional LP-01-V2 \\
\hline
\end{tabular}

\vspace{0.5cm}

% ----------------------------------------------------------------------------
% DIFFERENZIERUNG
% ----------------------------------------------------------------------------
\section*{Differenzierung}

\begin{tcolorbox}[colback=white, colframe=basis, title=\textcolor{white}{\textbf{[B] Basis}}]
\begin{itemize}[nosep]
    \item 8 Kernvokabeln (rezeptiv ausreichend)
    \item Bildkarten als permanente Stütze
    \item AB: Aufgaben 1--3 (Zuordnung, Artikel, einfache Sätze)
    \item Kettenübung: Darf ablesen
\end{itemize}
\end{tcolorbox}

\begin{tcolorbox}[colback=white, colframe=standard, title=\textcolor{white}{\textbf{[S] Standard}}]
\begin{itemize}[nosep]
    \item 12 Vokabeln (produktiv)
    \item AB: Aufgaben 1--5 komplett
    \item Freie Sätze mit hay bilden
    \item Partnerarbeit ohne Hilfe
\end{itemize}
\end{tcolorbox}

\begin{tcolorbox}[colback=white, colframe=erweiterung, title=\textcolor{white}{\textbf{[E] Erweiterung}}]
\begin{itemize}[nosep]
    \item 18+ Vokabeln (inkl. Selbstrecherche)
    \item AB: Alle Aufgaben + Bonus
    \item Mini-Text: 4--5 Sätze über Klassenraum
    \item Expertenrolle: Hilft anderen
\end{itemize}
\end{tcolorbox}

\vspace{0.5cm}

% ----------------------------------------------------------------------------
% HÄUFIGE FEHLER
% ----------------------------------------------------------------------------
\section*{Häufige Fehler und Reaktionen}

\begin{tabular}{|p{4cm}|p{4cm}|p{5cm}|}
\hline
\textbf{Fehler} & \textbf{Ursache} & \textbf{Reaktion} \\
\hline
*Hay el libro & Deutsch-Transfer & Kontrastieren: hay = neu/unbestimmt \\
\hline
*la mapa & Endung -a = feminin? & Merkspruch: \textit{Der Mapa ist ein Mann!} \\
\hline
Aussprache \textit{ll} & Kein dt. Äquivalent & Vorsprechen: wie deutsches \textit{j} \\
\hline
*Es hay una mesa & Englisch-Einfluss & Kein Subjekt vor hay! \\
\hline
\end{tabular}

\vspace{0.5cm}

% ----------------------------------------------------------------------------
% MATERIALCHECKLISTE
% ----------------------------------------------------------------------------
\section*{Materialcheckliste}

\begin{itemize}[nosep]
    \item[$\square$] AB-01-V2 (Klassensatz)
    \item[$\square$] Bildkarten Klassenzimmer (laminiert)
    \item[$\square$] Lehrwerk Qué pasa 1 (S. 36--38)
    \item[$\square$] Optional: LP-01-V2 als QR-Code
    \item[$\square$] Tafel/Whiteboard für Tafelbild
    \item[$\square$] Eigene Tasche für Demo
\end{itemize}

\vspace{0.5cm}

% ----------------------------------------------------------------------------
% TAFELBILD-SKIZZE
% ----------------------------------------------------------------------------
\section*{Tafelbild}

\begin{tcolorbox}[colback=white, colframe=grau]
\textbf{Objetos de clase -- hay}

\vspace{0.3cm}
\begin{tabular}{p{6cm}|p{6cm}}
\textbf{Vokabeln} & \textbf{hay = es gibt} \\
\hline
la mesa -- der Tisch & \\
la silla -- der Stuhl & En la clase \textbf{hay una} mesa. \\
la pizarra -- die Tafel & \textbf{Hay cinco} sillas. \\
el libro -- das Buch & \textbf{Hay muchos} libros. \\
el cuaderno -- das Heft & \\
... & ¿Qué hay en la clase? \\
& $\rightarrow$ Hay una pizarra. \\
\textbf{Achtung:} el mapa (m.) & \\
\end{tabular}

\vspace{0.3cm}
\textbf{NICHT:} *hay el libro
\end{tcolorbox}

\vspace{0.5cm}

% ----------------------------------------------------------------------------
% HAUSAUFGABE
% ----------------------------------------------------------------------------
\section*{Hausaufgabe}

\begin{itemize}[nosep]
    \item Vokabeln lernen (8 Kernwörter)
    \item Lehrwerk S. 37, Aufgabe 2
    \item Optional: LP-01-V2 (digital)
    \item [E]: 5 Sätze über eigenes Zimmer
\end{itemize}

\vspace{0.5cm}

% ----------------------------------------------------------------------------
% AB-LÖSUNGEN KURZFORM
% ----------------------------------------------------------------------------
\section*{AB-Lösungen (Kurzform)}

\textbf{Aufg. 1:} 1c, 2a, 3d, 4f, 5b, 6e, 7h, 8g

\textbf{Aufg. 2:} a) la, b) el, c) la, d) el, e) la, f) \textbf{el} (!)

\textbf{Aufg. 3:} Individuelle Sätze mit hay + un/una/Zahl

\textbf{Aufg. 4:} hay, Hay, hay, un, hay, un, (Zahl)

\textbf{Aufg. 5:} Individuelle Antworten

\textbf{Bonus:} 4 korrekte Sätze mit hay

\end{document}
